\documentclass{situationist-preprint}

\title{Laws in Biology}
\msid{SCI-08-a}
\msfamily{SCI-08}
\msdiscipline{Philosophy of Science}
\mstype{Principal manuscript}
\mskeywords{laws of nature; biology; evolutionary contingency; ceteris paribus; invariance; exceptions; Mendel's laws}

\begin{document}
\maketitle

\begin{abstract}
Are there laws in biology? We argue that there are no strict biological laws, but that there are robust generalizations with stable exceptions. A strict law is an exceptionless, non-accidental generalization stated in the vocabulary of the science. The argument against strict biological laws combines Beatty's evolutionary contingency thesis with the heterogeneity of biological kinds: every generalization stated in biological vocabulary concerns the products of evolution, which variation and selection could have made otherwise and have in fact made otherwise somewhere. The generalizations that escape this argument are mathematical truths about idealized populations or laws of physics and chemistry applied to organisms. But many biological generalizations, such as Mendel's law of segregation and the allometric scaling of metabolic rate, are invariant across a wide range of conditions, and their exceptions are produced by identifiable, recurrent mechanisms. Such generalizations do the explanatory and predictive work that laws were supposed to do, and they are the generalizations that a best system relativized to biological kinds contains.
\end{abstract}

\section{The problem}\label{s:problem}

Physics states laws: Newton's laws of motion, Maxwell's equations, the Schr\"odinger equation. Biology states generalizations: that genes segregate in meiosis, that metabolic rate scales with body mass to the three-quarters power, that island populations of large mammals tend to dwarfism. Whether any of the latter are laws has been disputed since \citet{Smart1963} argued that biology has no laws in the strict sense, only generalizations about the organisms that happen to exist on Earth, comparable to the generalizations of engineering about the particular devices it studies.

\citet{Beatty1995} gave the argument its standard form, the evolutionary contingency thesis: all distinctively biological generalizations describe contingent outcomes of evolution, and the conditions that produced them could have been otherwise; so they lack the necessity, or at least the non-accidental character, of laws. \citet{Rosenberg2001} drew the conclusion that biological explanation is historical rather than nomological. Several responses followed. \citet{Sober1997} and \citet{Elgin2003} argued that biology has laws that are a priori, mathematical models such as the Hardy--Weinberg principle and Fisher's sex-ratio argument, whose truth does not depend on evolutionary contingencies. \citet{Mitchell1997,Mitchell2000} argued that lawhood comes in degrees along dimensions of stability, strength and abstraction, and that biological generalizations are laws of lower stability. \citet{Woodward2001} argued that what matters for explanation is invariance, not lawhood, and that biological generalizations are invariant over limited ranges. \citet{Lange1995} argued that there are laws concerning particular species, stable relative to the questions biologists ask. \citet{Waters1998} distinguished the contingent distributions of traits among organisms from the causal regularities that hold wherever the relevant entities exist, and located the contingency in the former.

We defend a position close to Woodward's and Mitchell's, stated as a definite answer: there are no strict laws in biology, but there are robust generalizations with stable exceptions.

\section{Framework}\label{s:framework}

\begin{definition}[Strict law]\label{d:strict}
A \emph{strict law} of a science is a true, non-accidental generalization stated in the vocabulary of that science, without exceptions and without a ceteris paribus clause.
\end{definition}

\begin{definition}[Robustness]\label{d:robust}
A generalization is \emph{robust} if it is invariant, in the sense of \citet{Woodward2003}, over a wide range of interventions on the variables it relates and of background conditions: it would continue to hold were those variables or conditions changed within the range.
\end{definition}

\begin{definition}[Stable exception]\label{d:stable}
An exception to a generalization is \emph{stable} if it is produced by a mechanism that is identifiable and recurrent, so that the class of exceptions is itself the subject of projectible generalizations. An exception that occurs at random, or by a mechanism found only once, is \emph{unstable}.
\end{definition}

The notion of a stable exception separates a generalization with exceptions from a false generalization. A false generalization has exceptions that occur without pattern; a robust generalization with stable exceptions has exceptions that are themselves explained by further generalizations.

\begin{assumption}[Evolutionary origin]\label{a:origin}
The entities and properties picked out by distinctively biological vocabulary, such as gene, species, organ, meiosis and metabolism, are products of evolution by variation, selection and drift, and their properties depend on the history of the lineages that bear them.
\end{assumption}

\begin{assumption}[Variation]\label{a:variation}
For any heritable property of the members of a biological kind, variation can produce, and over the history of life on Earth has produced, lineages in which the property is altered.
\end{assumption}

Assumption~\ref{a:variation} is supported by the record. The genetic code, once proposed as a frozen accident shared by all life, has variants in mitochondria and in several groups of microorganisms. Meiosis, the basis of Mendelian segregation, is subverted by meiotic drive, in which some alleles are transmitted to more than half of the gametes. Many lineages of plants and animals reproduce asexually.

\section{The argument}\label{s:argument}

\begin{theorem}[No strict laws]\label{t:nolaws}
There are no strict laws of biology.
\end{theorem}

\begin{proof}
Let $G$ be a generalization stated in biological vocabulary, saying that members of a biological kind $K$ have property $Q$, or that a biological variable stands in a given relation to others. By Assumption~\ref{a:origin}, $K$ and $Q$ are products of evolution, and the properties of members of $K$ depend on the histories of their lineages. If $Q$ is heritable, then by Assumption~\ref{a:variation} variation can produce, and has produced, members of $K$ that lack it; then $G$ has exceptions, or holds only with a clause restricting it to normal members, and so is not strict. If $Q$ is not heritable but follows from the heritable properties of members of $K$ together with physical and chemical laws, then a variation in those heritable properties can alter $Q$, and the same holds. If $Q$ follows from physical and chemical laws together with no heritable property at all, then $G$ is a law of physics or chemistry applied to members of $K$, not a biological law. The remaining case is that $G$ follows from mathematics alone together with idealizing conditions, and then $G$, as a conditional whose antecedent states the conditions, is a mathematical truth, not an empirical law of biology; we consider this case in Objection~\ref{o:apriori}.
\end{proof}

The proof is the contingency argument of \citet{Beatty1995}, with the cases made explicit. It has a counterpart for kinds. MET-27 argues that natural kinds are clusters of co-instantiated genuine properties and that biological taxa lack exceptionless clusters; the absence of strict laws is the same fact viewed from the side of generalizations. A strict law about the members of a taxon would require a property they all share, without exception, and that is what variation denies.

\begin{proposition}[Robust generalizations]\label{p:robust}
Some biological generalizations are robust and have only stable exceptions.
\end{proposition}

\begin{proof}
We give two examples. Mendel's law of segregation says that the two alleles of a gene in a diploid organism are transmitted to the gametes with equal probability. It holds across the eukaryotes that reproduce sexually, across environments and across most genes, and it is invariant under interventions on which alleles an organism carries: crossing different parental genotypes yields the predicted ratios. Its exceptions are produced by identifiable mechanisms, meiotic drive, gametic selection and non-disjunction, each recurrent and each the subject of its own generalizations. Kleiber's law says that basal metabolic rate scales with body mass to roughly the three-quarters power. It holds across many orders of magnitude of mass and across many groups of animals, and its departures, in some taxa and size ranges, are systematic and have been modelled in terms of the geometry of distribution networks and the thermal environment.
\end{proof}

Robust generalizations with stable exceptions support the work laws were thought to do. They support counterfactuals within their range of invariance, so they answer the what-if questions that, on the account of SCI-03, constitute explanation. They support prediction, with the exceptions predictable in turn when the exception-producing mechanism is known. And they are projectible, since the kinds they concern are natural enough for the inferences biologists make with them.

The best-system account of laws defended in MET-12 explains the status of these generalizations. The best system for the fundamental vocabulary contains no biological generalizations, since they are not stated in that vocabulary. A best system relativized to biological kinds, in the manner of \citet{CohenCallender2009}, contains the generalizations that best combine simplicity and strength in summarizing the biological facts; these hold with the stability that a best summary confers and with the exceptions that the heterogeneity of the kinds introduces. Whether to call them biological laws is a terminological matter. They are not strict laws, and Theorem~\ref{t:nolaws} does not deny that they are laws in the relativized sense.

\section{Objections and replies}\label{s:objections}

\begin{objection}[A priori laws]\label{o:apriori}
The Hardy--Weinberg principle and Fisher's sex-ratio principle are biological laws without exceptions; they are a priori, so evolutionary contingency does not touch them \citep{Sober1997,Elgin2003}.
\end{objection}

\begin{reply}
We agree that they are exceptionless and a priori, and that they are important. But each is a conditional whose antecedent describes an idealized population: random mating, no selection, no mutation, no migration, infinite size. As such it is a theorem of probability theory, and its biological content lies in the claim that actual populations approximately satisfy the antecedent, which is contingent and has exceptions. The a priori conditional is not a strict empirical law of biology, and the empirical claim is not strict. The two together form a robust generalization with stable exceptions, the exceptions being populations under selection, inbreeding or migration, which is our thesis.
\end{reply}

\begin{objection}[Ceteris paribus laws]
Special-science laws are ceteris paribus laws, and ceteris paribus laws are laws \citep{Fodor1974}. Biology has them in abundance.
\end{objection}

\begin{reply}
A ceteris paribus law, on the most defensible analyses, is a generalization that holds when interfering factors are absent, where the interferers are independently specifiable. That is a robust generalization whose exceptions are stable in the sense of Definition~\ref{d:stable}; the independent specifiability of the interferers is what makes them stable. We therefore accept the objector's examples and deny only that they are strict. The well-known worry that ceteris paribus clauses make laws vacuous \citep{EarmanRoberts1999} is answered for biology by the requirement that exceptions be stable: a clause that excuses every exception would allow unstable ones.
\end{reply}

\begin{objection}[Physics too]
Physical laws also have exceptions: the law of gravitation is false of bodies that are charged, and every law of physics, stated as a regularity, is false of some situation \citep{Cartwright1983}. If biology has no strict laws, neither does physics.
\end{objection}

\begin{reply}
The argument confuses laws stated as regularities about total behavior with laws stated as component influences. The law of gravitation describes the gravitational component of the force on a body, and it holds without exception in that form. The fundamental laws of physics are, on the account of MET-12, the axioms of the best system for the fundamental vocabulary, and nothing in the argument of Theorem~\ref{t:nolaws} applies to them, since they do not concern products of evolution. The asymmetry is that biological generalizations fail in their own terms, through variation in the entities they concern, while the fundamental physical laws fail only when misread.
\end{reply}

\begin{objection}[Species-specific laws]
\citet{Lange1995} argues that there are laws concerning particular species, such as that members of a species of bird have a certain plumage, stable under the counterfactuals relevant to biological inquiry.
\end{objection}

\begin{reply}
Lange's laws are stable relative to a restricted range of counterfactual suppositions, those that a given field of biology treats as relevant. That is invariance over a range, which is robustness in the sense of Definition~\ref{d:robust}, and the exceptions, such as the moulting, the juvenile and the leucistic individual, are stable. On our account these are robust generalizations of limited range. Lange's insistence that their stability is relative to a field of inquiry is the same insight as the relativization of the best system to biological kinds.
\end{reply}

\section{Conclusion}\label{s:conclusion}

Every generalization stated in biological vocabulary concerns products of evolution, and variation has produced, somewhere, members of the relevant kinds that violate it. The exceptionless generalizations that biology uses are mathematical truths about idealized populations or laws of physics and chemistry applied to organisms. There are therefore no strict biological laws. But there are robust generalizations, invariant over wide ranges of interventions and conditions, whose exceptions are produced by identifiable, recurrent mechanisms; they explain, predict and support counterfactuals, and they are what a best system relativized to biological kinds contains.

\begin{thebibliography}{99}\small
\bibitem[Beatty(1995)]{Beatty1995} Beatty, J. (1995). The evolutionary contingency thesis. In G.~Wolters and J.~G. Lennox (eds.), \emph{Concepts, Theories, and Rationality in the Biological Sciences}. University of Pittsburgh Press.
\bibitem[Cartwright(1983)]{Cartwright1983} Cartwright, N. (1983). \emph{How the Laws of Physics Lie}. Oxford University Press.
\bibitem[Cohen and Callender(2009)]{CohenCallender2009} Cohen, J., and Callender, C. (2009). A better best system account of lawhood. \emph{Philosophical Studies}, 145, 1--34.
\bibitem[Earman and Roberts(1999)]{EarmanRoberts1999} Earman, J., and Roberts, J. (1999). Ceteris paribus, there is no problem of provisos. \emph{Synthese}, 118, 439--478.
\bibitem[Elgin(2003)]{Elgin2003} Elgin, M. (2003). Biology and a priori laws. \emph{Philosophy of Science}, 70, 1380--1389.
\bibitem[Fodor(1974)]{Fodor1974} Fodor, J.~A. (1974). Special sciences (or: the disunity of science as a working hypothesis). \emph{Synthese}, 28, 97--115.
\bibitem[Lange(1995)]{Lange1995} Lange, M. (1995). Are there natural laws concerning particular biological species? \emph{Journal of Philosophy}, 92, 430--451.
\bibitem[Mitchell(1997)]{Mitchell1997} Mitchell, S.~D. (1997). Pragmatic laws. \emph{Philosophy of Science}, 64, S468--S479.
\bibitem[Mitchell(2000)]{Mitchell2000} Mitchell, S.~D. (2000). Dimensions of scientific law. \emph{Philosophy of Science}, 67, 242--265.
\bibitem[Rosenberg(2001)]{Rosenberg2001} Rosenberg, A. (2001). How is biological explanation possible? \emph{British Journal for the Philosophy of Science}, 52, 735--760.
\bibitem[Smart(1963)]{Smart1963} Smart, J.~J.~C. (1963). \emph{Philosophy and Scientific Realism}. Routledge and Kegan Paul.
\bibitem[Sober(1997)]{Sober1997} Sober, E. (1997). Two outbreaks of lawlessness in recent philosophy of biology. \emph{Philosophy of Science}, 64, S458--S467.
\bibitem[Waters(1998)]{Waters1998} Waters, C.~K. (1998). Causal regularities in the biological world of contingent distributions. \emph{Biology and Philosophy}, 13, 5--36.
\bibitem[Woodward(2001)]{Woodward2001} Woodward, J. (2001). Law and explanation in biology: invariance is the kind of stability that matters. \emph{Philosophy of Science}, 68, 1--20.
\bibitem[Woodward(2003)]{Woodward2003} Woodward, J. (2003). \emph{Making Things Happen: A Theory of Causal Explanation}. Oxford University Press.
\end{thebibliography}

\end{document}
